\section{The universal quartic bound}\label{sec:upper}

For completeness, we prove the upper bound used in
Theorem~\ref{thm:main}. The argument is the one of Aaronson, Ben-David,
Kothari, Rao, and Tal~\cite{ABKRT21}, combining Huang's signed-cube
method~\cite{Huang19} with the adversary bound and the deterministic
degree--block-sensitivity inequality. No part of the construction is
used in this appendix.

For a total Boolean function $f$, let $\D(f)$ be its exact deterministic
query complexity and $\deg(f)$ its degree as a real multilinear polynomial.
The degree is unchanged by replacing input or output bits by signs.
A nonempty set $B\subseteq[n]$ is a \emph{sensitive block at $x$} if
flipping every bit in $B$ changes $f(x)$. The block sensitivity $\bs(f)$
is the maximum, over $x$, of the number of pairwise disjoint sensitive
blocks at $x$.

\begin{theorem}[\cite{ABKRT21}]\label{thm:quartic-upper}
Every total Boolean function satisfies
\[
 \R(f)\le\D(f)=O((1+\Q(f))^4).
\]
\end{theorem}

We establish the three ingredients separately. All matrix norms below
are operator norms for the Euclidean norm.

\begin{lemma}[Nonnegative spectral adversary bound]
\label{lem:adversary}
Let $f$ be a nonconstant total Boolean function. Let $\Gamma$ be a nonzero
real symmetric matrix, indexed by inputs, with nonnegative entries and
$\Gamma_{xy}=0$ whenever $f(x)=f(y)$. Let
\[
 (\Gamma_i)_{xy}=\Gamma_{xy}\,\mathbf1_{\{x_i\ne y_i\}}.
\]
Then, for an absolute constant $c>0$,
\[
 \Q(f)\ge c\frac{\norm{\Gamma}}{\max_i\norm{\Gamma_i}}.
\]
\end{lemma}

\begin{proof}
This is the spectral adversary argument of
Barnum, Saks, and Szegedy~\cite{BSS03}; see also
\v{S}palek and Szegedy~\cite{SpalekSzegedy06}.
Purify all randomness and defer measurements. A bounded-query algorithm
can be represented with a fixed number of queries, with at most a
constant-factor overhead: branches that have finished can perform
ineffective queries, and controlled queries can be simulated with two
bit queries. Let $\lvert\psi_x^t\rangle$ be its normalized pure state
on input $x$ after $t$ queries.

By the Perron--Frobenius theorem, $\Gamma$ has a nonnegative unit
eigenvector $v$ of eigenvalue $\norm{\Gamma}$. Track
\[
 W_t=\sum_{x,y}v_xv_y\Gamma_{xy}
            \langle\psi_x^t\mid\psi_y^t\rangle.
\]
Initially the states are input-independent, so
$W_0=v^\top\Gamma v=\norm{\Gamma}$. Input-independent unitaries
preserve every overlap.

Immediately before a query, write the normalized state as
\[
 \lvert\psi_x\rangle=\sum_i\lvert i\rangle\lvert\phi_{x,i}\rangle,
 \qquad \sum_i\norm{\phi_{x,i}}^2=1.
\]
The bit oracle acts identically
on inputs $x,y$ at that index unless $x_i\ne y_i$. When they differ,
the overlap changes in magnitude by at most
$2\norm{\phi_{x,i}}\norm{\phi_{y,i}}$, since the difference between
the relevant unitary and the identity has norm at most two. Put
$u_i(x)=v_x\norm{\phi_{x,i}}$. Nonnegativity of the weights gives
\begin{align*}
 \abs{W_{t+1}-W_t}
 &\le2\sum_i u_i^\top\Gamma_i u_i\\
 &\le2\max_i\norm{\Gamma_i}\sum_i\norm{u_i}^2
  =2\max_i\norm{\Gamma_i}.
\end{align*}
The last equality follows from normalization of each query state and
$\sum_xv_x^2=1$.

At the end, measure the output bit. For opposite-output inputs, one
output-one probability is at most $1/3$, and the other is at least
$2/3$. Splitting the two states by that output and applying
Cauchy--Schwarz bounds their overlap in absolute value by
$2\sqrt2/3$. Hence the final value satisfies
\[
 \abs{W_{\mathrm{final}}}
 \le\frac{2\sqrt2}{3}\sum_{x,y}v_xv_y\Gamma_{xy}
 =\frac{2\sqrt2}{3}\norm{\Gamma}.
\]
The required total change in $W$ divided by the per-query bound proves
the assertion.
\end{proof}

Let $S_f$ be the adjacency matrix of the \emph{sensitivity graph}: its
vertices are $\bits^n$, and two vertices are adjacent if they differ
in exactly one bit and have different function values. Write
$\lambda(f)=\norm{S_f}$.

\begin{lemma}\label{lem:bs-lambda}
For every nonconstant total Boolean function,
\[
 \bs(f)=O(\Q(f)^2),\qquad
 \lambda(f)=O(\Q(f)).
\]
\end{lemma}

\begin{proof}
Choose an input $x$ with $b=\bs(f)$ disjoint sensitive blocks. Take
$\Gamma$ to be the adjacency matrix of the star whose center is $x$
and whose leaves are the $b$ inputs obtained by flipping those blocks,
with zero entries elsewhere. Its norm is $\sqrt b$. Each coordinate
belongs to at most one block, so $\Gamma_i$ has at most one undirected
edge and norm at most one. Lemma~\ref{lem:adversary} gives
$\Q(f)=\Omega(\sqrt b)$. This recovers the usual block-sensitivity
lower bound for quantum queries~\cite{BBCMW01}.

For the second assertion, use $\Gamma=S_f$. The edges in any one
coordinate form a partial matching, whose adjacency matrix has norm
at most one. Another application of Lemma~\ref{lem:adversary} gives
$\Q(f)=\Omega(\lambda(f))$, as in~\cite{ABKRT21}.
\end{proof}

\begin{lemma}[\cite{ABKRT21}, using \cite{Huang19}]
\label{lem:degree-lambda}
Every nonconstant total Boolean function satisfies
$\deg(f)\le\lambda(f)^2$.
\end{lemma}

\begin{proof}
Write the function with sign inputs and sign outputs, and set
$d=\deg(f)\ge1$. Choose a monomial supported on a set of $d$ variables
with nonzero coefficient, and fix all other variables to any signs.
Its coefficient survives: any different monomial that could contribute
to that coefficient would strictly contain its support and have degree
larger than $d$. The restricted sign function $g$ on the $d$-cube
therefore has nonzero correlation with parity. For some sign
$\sigma\in\{-1,1\}$, the set
\[
 U=\{x\in\{-1,1\}^d:g(x)=\sigma x_1\cdots x_d\}
\]
has size greater than $2^{d-1}$. Every cube edge with both endpoints in
$U$ changes $g$, since parity changes across each edge.

Define a signed adjacency matrix of the cube recursively by
\[
 H'_1=\begin{pmatrix}0&1\\1&0\end{pmatrix},
 \qquad
 H'_{j+1}=\begin{pmatrix}H'_j&I\\I&-H'_j\end{pmatrix}.
\]
Its entrywise absolute value is the ordinary cube adjacency matrix.
Induction gives $(H'_j)^2=jI$, and its trace is zero. Thus $H'_d$ has
eigenvalues $\pm\sqrt d$, each with multiplicity $2^{d-1}$. The
$\sqrt d$-eigenspace intersects the space of vectors supported on $U$
nontrivially, because their dimensions sum to more than $2^d$.
Choose a nonzero real vector $v$ in this intersection, and put
$u=\abs v$ entrywise. If $A_U$ is the unsigned adjacency induced on
$U$, then
\[
 u^\top A_Uu\ge v^\top H'_dv=\sqrt d\,\norm v^2
             =\sqrt d\,\norm u^2.
\]
All edges of $A_U$ are sensitivity edges of the restriction, and hence
of the original function. Extending $u$ by zero to the full input
cube gives $\lambda(f)\ge\sqrt d$, as required.
\end{proof}

\begin{lemma}[\cite{Midrijanis04}]\label{lem:degree-bs}
Every total Boolean function satisfies
$\D(f)\le\bs(f)\deg(f)$.
\end{lemma}

\begin{proof}
Consider the following deterministic algorithm. If the function
restricted by the queries so far is nonconstant, select a monomial
of maximum degree in that restricted polynomial and query all its
variables. Continue until the restriction is constant.

Fix the actual full input $x$. In a given round, let $S$ be the
selected monomial's support. Fix all other still-unqueried variables
to their values in $x$. The coefficient on $S$ remains nonzero, since
no monomial in the current restriction strictly contains $S$.
The resulting function on $S$ is nonconstant. Some assignment to $S$
therefore differs in value from the actual assignment $x|_S$. The
nonempty subset of coordinates on which these two assignments differ
is a sensitive block at the original input $x$.

Different rounds query disjoint sets, so these sensitive blocks are
disjoint. There are at most $\bs(f)$ rounds. Restrictions do not
increase degree, and each round queries at most $\deg(f)$ variables.
The stated deterministic bound follows.
\end{proof}

\begin{proof}[Proof of Theorem~\ref{thm:quartic-upper}]
For nonconstant $f$, Lemmas~\ref{lem:bs-lambda} and
\ref{lem:degree-lambda} give both $\bs(f)=O(\Q(f)^2)$ and
$\deg(f)=O(\Q(f)^2)$. Lemma~\ref{lem:degree-bs} then gives
$\D(f)=O(\Q(f)^4)$. Since a deterministic algorithm is also randomized,
$\R(f)\le\D(f)$. Constant functions require no queries.
\end{proof}
